% Slug: four-of-diamonds
% Description: ?
% Author: Cisco
% Story: Cisco
% Test Data: Cisco
% Testers: ?
% Editorial: ?
% TempSlug: four-of-diamonds

% Please don't modify or delete TempSlug

\problem{Six Seven}
Alice was tasked with implementing Dijkstra's shunting yard, an algorithm for parsing arithmetic expressions, respecting correct operator precedence (i.e. MDAS).  However, she submitted a meme instead, which annoyed her teacher.
\begin{itemize}
    \item The \texttt{+} operation is implemented as \textbf{string concatenation} instead of addition, \newline e.g. \texttt{6 + 7} evaluates to \texttt{67}
    \item The \texttt{*} operation is implemented as \textbf{string duplication} instead of multiplication, \newline e.g. \texttt{6 * 7} evaluates to \texttt{6666666}
    \item The \texttt{-} operation is implemented as normal, except negative numbers were not implemented, so if any operation would result in one at any intermediate step \newline (e.g. \texttt{6 - 7}) then the computer explodes.
    \item The \texttt{/} operation is implemented as normal, except non-integer fractions were not implemented, so if any operation would result in one at any intermediate step \newline (e.g. \texttt{6 / 7}) then the computer explodes.
    \item Finally, the \textbf{only} accepted integer literals are \texttt{6} and \texttt{7}.
\end{itemize}
Otherwise, her program correctly follows the MDAS convention.
\begin{itemize}
    \item It recursively evaluates parenthesized expressions first, then
    \item resolves \texttt{*} and \texttt{/} operations from left to right, then
    \item resolves \texttt{+} and \texttt{-} operations from left to right.
\end{itemize}
Alice's program ignores spaces between tokens.

For example, \texttt{7 - 6 + 6} equals \texttt{16}, evaluated as follows:
\begin{align*}
    \texttt{7 - 6 + 6}
    &= \texttt{{\color{red}1} + 6} \\
    &= \texttt{{\color{red}16}}.
\end{align*}

For another example, \texttt{(7+7) - (6+7 - 7)/6} equals \texttt{67}, evaluated as follows:
\begin{align*}
    \texttt{(7+7) - (6+7 - 7)/6}
    &= \texttt{{\color{red}77} - (6+7 - 7)/6}, \\
    &= \texttt{77 - ({\color{red}67} - 7)/6}, \\
    &= \texttt{77 - {\color{red}60}/6}, \\
    &= \texttt{77 - \color{red}10} \\
    &= \texttt{\color{red}67}.
\end{align*}

For one final example, \texttt{6*6 / (6 + (6 + 6)) / 7} equals \texttt{143}, evaluated as follows:
\begin{align*}
    \texttt{6*6 / (6 + (6 + 6)) / 7}
    &= \texttt{6*6 / (6 + {\color{red}66}) / 7} \\
    &= \texttt{6*6 / {\color{red}666} / 7} \\
    &= \texttt{{\color{red}666666} / 666 / 7} \\
    &= \texttt{{\color{red}1001} / 7} \\
    &= \texttt{{\color{red}143}}.
\end{align*}

Given a positive integer $n$, your goal is to construct any arithmetic expression (using parentheses, \texttt{6} and \texttt{7}, and the \texttt{*/+-} operations) which evaluates to $n$ using Alice's meme program.  \textit{For completeness}, here are the other limitations of Alice's program:
\begin{itemize}
    \item If the expression contains more than $1000$ characters (including spaces), \newline the computer blows up.
    \item If at any intermediate step an integer with a value greater than $2^{64}$ is produced, \newline the computer blows up.
    \item For \texttt{+} and \texttt{*}, leading zeros are removed at each intermediate step, \newline e.g. \texttt{0 + 0} and \texttt{0 * 5} both evaluate to \texttt{0} (just one \texttt{0})
    \item Finally, \texttt{n * 0} causes the computer to blow up, regardless of the left operand.
\end{itemize}

\subsection*{Input Format}

Input consists of a single positive integer $n$.


\subsection*{Output Format}

Output a single line containing \textit{any} valid expression which evaluates to $n$.

We give the following formal definition of what exactly constitutes an \textit{expression}.

An \textit{expression} is a sequence of one or more sub-expressions, separated by operators.  \newline
A \textit{sub-expression} is one of the following:
\begin{itemize}
\item an integer literal (which must be \texttt{6} or \texttt{7}); or
\item another expression, enclosed in parentheses \texttt{()}
\end{itemize}

Each operator must be one of \texttt{+} or \texttt{-} or \texttt{*} or \texttt{/}

Spaces between tokens are ignored, so add as many as you like (within the char. count).

Here are some more examples of valid expressions:
\begin{center}
    \texttt{(7) - ((6)) + (((7)))}
    
    \texttt{(6 + 7 * 6)}

    \texttt{(6 / 6) * (7)}

    \texttt{(((((((6)))))))}

    \texttt{7}
\end{center}
These evaluate to \texttt{17} and \texttt{6777777} and \texttt{1111111} and \texttt{6} and \texttt{7}, respectively.

\subsection*{Constraints}

\textit{Your program will only be tested on data that satisfies the following constraints.}

\constraints{
    $1 \leq n \leq 10^{12}$ \\
}

\textit{There is no need to validate in your program that these constraints hold.}

% \newpage

\subsection*{Sample I/O}

\begin{sampleIO}
\begin{verbatim}
16
\end{verbatim}
&
\begin{verbatim}
7 - 6 + 6
\end{verbatim}
\end{sampleIO}

\begin{sampleIO}
\begin{verbatim}
67
\end{verbatim}
&
\begin{verbatim}
(7+7) - (6+7 - 7)/6
\end{verbatim}
\end{sampleIO}


\begin{sampleIO}
\begin{verbatim}
143
\end{verbatim}
&
\begin{verbatim}
6*6 / (6 + (6 + 6)) / 7
\end{verbatim}
\end{sampleIO}


\begin{sampleIO}
\begin{verbatim}
17
\end{verbatim}
&
\begin{verbatim}
(7) - ((6)) + (((7)))
\end{verbatim}
\end{sampleIO}


\begin{sampleIO}
\begin{verbatim}
6777777
\end{verbatim}
&
\begin{verbatim}
(6 + 7 * 6)
\end{verbatim}
\end{sampleIO}


\begin{sampleIO}
\begin{verbatim}
1111111
\end{verbatim}
&
\begin{verbatim}
(6 / 6) * (7)
\end{verbatim}
\end{sampleIO}



\begin{sampleIO}
\begin{verbatim}
6
\end{verbatim}
&
\begin{verbatim}
(((((((6)))))))
\end{verbatim}
\end{sampleIO}



\begin{sampleIO}
\begin{verbatim}
7
\end{verbatim}
&
\begin{verbatim}
7
\end{verbatim}
\end{sampleIO}



% \subsection*{Explanation}
% TODO